(1) 对球对称压强 \(p'(r,t)=P(r)e^{-i\omega t}\)，球坐标 Laplace 算符为 \(\nabla^2P=r^{-2}(r^2P')'\)，故
\[
\frac1{r^2}(r^2P')'+k^2P=0,\qquad k=\omega/c.
\]
通解为球贝塞尔函数的组合
\[
P(r)=A j_0(kr)+B y_0(kr)
=A\frac{\sin kr}{kr}-B\frac{\cos kr}{kr}.
\]
由于 \(y_0(kr)\propto1/r\) 在球心发散，必须取 \(B=0\)；而 \(j_0(0)=1\) 有限且导数在原点为零。

(2) 流体线性动量方程 \(\rho_0\partial_t v_r=-\partial_r p'\)。对 \(\omega>0\) 的谐波振动，刚性壁 \(v_r(a,t)=0\) 等价于 \(P'(a)=0\)，即
\[
j_0'(ka)=0\quad\Longleftrightarrow\quad
ka\cos(ka)-\sin(ka)=0.
\]
对正根令 \(z=ka\)，则满足 \(\tan z=z\)，不包括 \(z=0\) 的静态零模。第一个非零正根为 \(z_1\simeq4.493409458\)，所以最低球对称声学振荡频率
\[
\omega_{1,\rm radial}=\frac{c z_1}{a}\simeq4.493409458\,\frac ca.
\]
这仅是球对称径向支的最低正频，不代表所有非球对称角模的最低正频。

(3) 两个不同特征值 \(k_m^2\ne k_n^2\) 对应的径向方程交叉相乘相减并在 \(r\in(0,a)\) 积分，利用正则球心条件和 Neumann 壁条件，得到
\[
(k_m^2-k_n^2)\int_0^a r^2P_m(r)P_n(r)\,dr=0,
\]
故径向模在 \(r^2dr\) 权重下正交。声压与速度的线性能量密度为
\[
u=\frac{p'^2}{2\rho_0c^2}+\frac{\rho_0|\boldsymbol v|^2}{2}.
\]
结合连续方程 \(\partial_tp'=-\rho_0c^2\nabla\cdot\boldsymbol v\) 和动量方程，可得局部能量平衡
\[
\partial_tu+\nabla\cdot(p'\boldsymbol v)=0.
\]
刚性壁 \(\boldsymbol v\cdot\hat n=0\)，故对空腔积分的总声能随时间守恒。

(4) \(k=0\) 时径向正则解为常数 \(P(r)=P_0\)，对应 \(\omega=0\) 的空间均匀静压微扰，而非周期性传播声波；因此频率谱中的零模不能算作第一声学振动。当半径 \(a\mapsto\eta a\) 且介质声速不变时，本征无量纲根 \(z_n\) 不变，全部球对称非零频率 \(\omega_n=c z_n/a\) 按 \(1/\eta\) 缩放，符合波动方程的尺度规律。
